\chapter{Planejamento experimental complementar}
\label{ap:planejamento_experimental}

Este apêndice registra extensões de análise que podem fortalecer a dissertação em etapas posteriores, mas que exigem nova execução do arquivo \texttt{simulacao.R} ou alteração do próprio script. O objetivo, nesta fase de qualificação, é separar com clareza o que já foi incorporado ao texto com base nos resultados existentes daquilo que demandaria esforço computacional ou modelagem adicional.

\section{Experimentos que exigem nova execução com a estrutura atual do script}

O arquivo \texttt{simulacao.R} já implementa sete \emph{setups} principais e pode ser reexecutado, sem alteração estrutural, para atualizar resultados quando houver necessidade de revisar premissas paramétricas. Entre os casos mais imediatos, destacam-se os seguintes:

\begin{enumerate}
    \item atualização dos parâmetros econômicos de referência, como preço do \emph{bunker fuel}, preço do biodiesel por rota e taxa de câmbio EUR/USD, para refletir nova janela temporal de mercado;
    \item atualização do sinal regulatório europeu, sobretudo dos valores de $\gamma^{ref}$ e $\gamma^{comp}$, caso se deseje aproximar o modelo de nova faixa observada no \gls{EUETS};
    \item atualização dos fatores de emissão \gls{WTW} das rotas de biodiesel, caso novas estimativas bibliográficas sejam adotadas para \gls{UCO}, sebo bovino, soja ou \gls{TCO};
    \item reexecução dos gráficos e tabelas de sensibilidade econômica, ambiental, regulatória, operacional e específica de \gls{TCO}, a partir de qualquer revisão dos parâmetros-base já definidos no script.
\end{enumerate}

\section{Experimentos que exigem alteração do \texttt{simulacao.R}}

Algumas extensões relevantes para uma etapa posterior do trabalho já podem ser delineadas, mas demandam modificação explícita do código atual.

\subsection{Desacoplamento entre os sinais regulatórios}

No \emph{Setup} E6, a trajetória regulatória ampliada é implementada com aumento simultâneo de $\gamma^{ref}$ e $\gamma^{comp}$. Uma extensão metodologicamente útil consiste em desacoplar esses dois componentes, permitindo testar cenários em que o preço de referência do carbono e o custo efetivo de \emph{compliance} evoluem em velocidades diferentes. Isso exigiria:

\begin{enumerate}
    \item redefinir a tabela \texttt{policy\_tbl} com combinações independentes entre os dois sinais;
    \item reexecutar o cálculo dos cenários E6;
    \item gerar novas figuras de fronteira regulatória e novas tabelas-resumo.
\end{enumerate}

\subsection{Refinamento do domínio de mistura}

O modelo atual avalia $\alpha$ em um conjunto discreto de valores. Caso se deseje investigar com maior precisão a transição entre inviabilidade e viabilidade, especialmente no intervalo entre $\alpha=0{,}24$ e $\alpha=0{,}30$, será necessário:

\begin{enumerate}
    \item redefinir o vetor \texttt{alpha\_grid} com malha mais fina;
    \item reexecutar os \emph{setups} E1, E4 e E6;
    \item regenerar os gráficos de preço de carbono de equilíbrio, fronteira regulatória e \emph{trade-off}.
\end{enumerate}

\subsection{Expansão para múltiplas rotas marítimas}

O script atual está parametrizado com foco na rota Santos--Roterdã. A inclusão de outras rotas exigiria alteração do código para leitura ou definição de novos vetores operacionais $(D,V,c)$, bem como a organização de uma camada adicional de comparação entre trajetos. Esse passo é particularmente relevante para estudos posteriores voltados à robustez externa dos resultados.

\subsection{Modelagem estocástica}

O modelo implementado é determinístico e paramétrico. Uma extensão natural consiste em introduzir incerteza explícita, por exemplo por simulação de Monte Carlo para preços, fatores de emissão, parâmetros operacionais e câmbio. Essa ampliação exigiria:

\begin{enumerate}
    \item definição de distribuições probabilísticas para os parâmetros de entrada;
    \item criação de nova rotina de simulação no \texttt{simulacao.R};
    \item redefinição das saídas para incluir intervalos, distribuições e probabilidades de viabilidade.
\end{enumerate}

\subsection{Custos adicionais de adaptação e logística}

O modelo atual representa a viagem a partir de custo de combustível e custo de \emph{compliance} regulatório. Em uma fase posterior, pode ser útil adicionar custos associados a armazenagem, \emph{bunkering}, segregação, adaptação operacional ou logística portuária. Essa extensão exigiria novas equações de custo e nova parametrização empírica, além de rerun completo dos experimentos centrais.

\section{Uso recomendado deste planejamento}

Para a fase de qualificação, o texto principal da dissertação deve priorizar os resultados já efetivamente gerados e discutidos. As extensões registradas neste apêndice funcionam como agenda metodológica explícita para aprofundamento posterior, sem confundir o leitor quanto ao que já foi executado e ao que ainda depende de nova rodada computacional ou de modificação do modelo.
